\chapter{Design and Trade-Off Analysis}
\label{ch:approach}

\section{The central trade-off: latency against interruption}
\label{sec:tradeoff}

Every decision in this chapter serves the trade-off introduced in
Chapter~\ref{ch:background}, which I state formally here. At each micro-pause, ending the
turn too early cuts the patient off, an interruption of cost $C_{\mathrm{cut}}$, while
waiting too long delays the reply, latency of cost $C_{\mathrm{lat}}$. Improving one
degrades the other, and no setting escapes both. Designing the system therefore comes down to deciding, at every instant, which side of that
trade-off to sit on, and deciding it for a stated reason. Each section below sets out one decision and the
trade-off it resolves.

\section{Which data trains, and which data validates}
\label{sec:datastrategy}

The project started from the assumption that Tennor's production transcripts would be the
training data and that the evaluation would be built around them. I reversed that in late
July, and because the reversal shaped everything downstream it belongs here rather than in
the implementation chapter.

Licensing forced it, and so did the labels. The eot-bench harness is Apache-2.0 and its
dataset CC-BY-4.0, both usable for research and derivative work; the production data is
personal data that cannot leave Tennor's infrastructure at all
(Appendix~\ref{app:ethics}), so anything built on it alone would be a result nobody else
could check. The two also differ in what they can support: the benchmark carries
human-verified turn boundaries, while the production set is labelled by weak supervision,
from whether the caller resumed speaking, which is cheap and scalable but noisy by
construction.

So I inverted the roles. The public benchmark carries the headline experiments (fusion,
calibration, cost), and the production data becomes an in-domain validation set, asked
only the narrower question of whether the method transfers.

The cost of that inversion is real and worth naming. The headline result is obtained on
generic French rather than on medical French, so the domain that motivates the project is
not the domain the main claim is measured on. And a weakly labelled validation set can only
return a weak answer, which is precisely what Chapter~\ref{ch:results} reports. What it
buys is that every headline number in this dissertation can be reproduced by anyone with a
Hugging Face token, and none of them rests on data that nobody outside Tennor can see.

\section{Late fusion rather than a jointly trained model}
\label{sec:latefusion}

The first modelling decision was how to combine the acoustic and the lexical signal. One
option is a single model trained jointly on audio and text. It can be more accurate,
because it learns the interactions between the two modalities, but it is heavier to run, it
demands aligned joint training data, and it locks both signals inside one box that cannot
be replaced piecemeal. The alternative, and the one I took, is late fusion: each signal
comes from an independent off-the-shelf model, and only their combination is learned.

I may lose some accuracy that way. What I gain is modularity, since each component can be
swapped on its own, along with speed and simplicity, which matches the constraint the
project set itself of composing existing models without retraining them.

\section{The lexical model: frozen embeddings rather than fine-tuning}
\label{sec:frozen}

The lexical signal can be obtained in two ways. CamemBERT can be fine-tuned on Tennor
medical French, which produces a \emph{medical-fitted} model closely adapted to that domain
but out of distribution on generic French, where it degrades the decision rather than
improving it. Section~\ref{sec:fusionresults} measures this. Or CamemBERT can be left frozen and used purely as a
feature extractor, with a small classifier on top, which generalises better and cannot
overfit the way a fine-tuned encoder can.

The two are named for the data they were fitted on in the chapters that follow:
\emph{benchmark-fitted} is the frozen encoder under a head fitted on the benchmark, and
\emph{medical-fitted} is the encoder fine-tuned on medical French. I kept both, for
different paths. The frozen encoder with a head fitted on the benchmark,
\emph{benchmark-fitted} in the tables of Chapter~\ref{ch:results}, serves the benchmark
evaluation. The medical-fitted model serves the in-domain path and the live demonstration,
where medical French is what it will actually see.

\section{A decision that never blocks on the transcript}
\label{sec:nonblocking}

One constraint governs the whole pipeline: the decision has to stay fast, while streaming
transcription can and does fall behind. I could wait for the most complete transcript available, paying added latency and risking a
stall, or read whatever partial transcript exists at the instant of the decision, as a
snapshot, accepting that it may lag the audio.

I took the second. The decision never waits on the recogniser: the transcript is read as
whatever snapshot exists at that instant. The two model calls are issued concurrently and
both are awaited, so the added latency is the slower of the two rather than their sum. Added latency is bounded by construction,
and the price is a text that is sometimes slightly incomplete. That is acceptable, because the
acoustic signal is always available and the fusion degrades gracefully when the lexical
branch has little to say. This is the design answer to the ASR dependency set out in
Section~\ref{sec:related}.

\section{The form of the fusion}
\label{sec:fusionform}

The fusion takes two forms, and the difference between them runs through the rest of this
dissertation. A rule-based form applies a bilateral veto and then a fixed
average; it needs no data to fit, so it was the committed deliverable and it is what the
system runs (Section~\ref{sec:fusionimpl}). The two operations are sequential rather than
combined: the veto is a gate, and when either signal falls below its threshold the decision
is to wait and the average is never computed. Only when both signals clear the gate are
they averaged, and the result is then mapped onto how much trailing silence to demand. A learned form fits the combination, and it is
the one every result in Chapter~\ref{ch:results} is measured on. The rule form is described
with the implementation because it is an engineering artefact; the learned form is set out
here because its shape is a design decision.

In the learned form the two probabilities are combined by logistic regression on
$p_{\mathrm{ac}}$, $p_{\mathrm{lex}}$ and their product. With weights $w_0 \dots w_3$,
\begin{equation}
z = w_0 + w_1 p_{\mathrm{ac}} + w_2 p_{\mathrm{lex}} + w_3\, p_{\mathrm{ac}} p_{\mathrm{lex}},
\qquad
p_{\mathrm{raw}} = \sigma(z) = \frac{1}{1 + e^{-z}} .
\label{eq:fusion}
\end{equation}
Here $p_{\mathrm{ac}} \in [0,1]$ is the probability of end of turn produced by the acoustic
model and $p_{\mathrm{lex}} \in [0,1]$ the one produced by the lexical model, both read at
the same micro-pause; $z$ is the linear score they are combined into; $\sigma$ is the
logistic function; and $p_{\mathrm{raw}} \in [0,1]$ is the fused probability before any
calibration. There are four weights because the model has four terms and no more: an
intercept $w_0$, one weight per input signal, and $w_3$ on their product.

This form contains the rule-based one as a special case. The fixed average
$0.5\,p_{\mathrm{ac}} + 0.5\,p_{\mathrm{lex}}$ is Equation~\ref{eq:fusion} with
$w_0 = 0$, $w_1 = w_2 = 0.5$ and $w_3 = 0$: the same shape, with the weights set by hand
instead of fitted, and with the interaction term switched off. What fitting them adds is
visible in the values Section~\ref{sec:fusionresults} reports: the lexical branch receives
roughly $2.2$ times the weight of the acoustic one, so the equal split the rule form
assumes is not what the data supports. The veto and the silence ladder run the other way:
they are gating structure that Equation~\ref{eq:fusion} has no counterpart for. The interaction
term $p_{\mathrm{ac}} p_{\mathrm{lex}}$ lets the model treat the case where both signals
agree as different from the sum of their separate contributions, which a purely additive
form cannot represent. 

I preferred logistic regression to richer combinations because it is light, its weights stay interpretable, so one can read off which of the
two signals carries more of the decision, and its output is a probability that can then
be calibrated. The fusion itself is small enough not to memorise a dataset of roughly a thousand spans,
though the lexical head above it is not: a logistic layer on 768-dimensional frozen
embeddings carries 769 parameters against some eight hundred training rows per fold, which
is why every lexical and fusion figure in Chapter~\ref{ch:results} is reported
out-of-fold. The weights are fitted by cross-validation with folds split by
call, so that no call appears in both training and test.

\section{Calibrating before deciding}
\label{sec:calibdesign}

A structural decision is to treat the probability the system produces as a quantity that
must be calibrated before it is used. The reason is direct. The final decision compares
that probability against a threshold derived from a cost, and the comparison only means
anything if the probability corresponds to a real frequency. As Section~\ref{sec:calibration-lit}
sets out, a poorly calibrated model can rank cases correctly and still be systematically
overconfident, which distorts any threshold placed on it.

Calibration is done by isotonic regression, fitting a monotonic $g$ on a held-out split to
minimise $\sum_i (g(p_i) - y_i)^2$. Because $g$ is monotonic it leaves the ordering of
predictions untouched, so the AUC does not move, and only the probability values are corrected. Quality is measured by the expected calibration error,
\begin{equation}
\mathrm{ECE} = \sum_{b} \frac{n_b}{N} \, \bigl| \mathrm{acc}_b - \mathrm{conf}_b \bigr| ,
\label{eq:ece}
\end{equation}
The predictions are sorted into bins $b$ by the probability they carry. For each bin,
$n_b$ is the number of predictions it holds, $N$ the total number of predictions,
$\mathrm{conf}_b$ the mean probability claimed within the bin, and $\mathrm{acc}_b$ the
proportion of those cases that turned out positive. The ECE is therefore the average gap
between claimed confidence and observed frequency, each bin weighted by how many predictions
fall into it. The cost of the approach is that it consumes held-out data and can
overfit on small sets, which I limit with a held-out split rather than reusing the fitting
data.

\section{Deciding from a cost rather than a fixed timeout}
\label{sec:costdesign}

The last decision is the stopping rule. The existing system stops after a fixed silence,
500\,ms say: simple, but arbitrary, and identical in every situation. Instead I choose the
operating point from the trade-off of Section~\ref{sec:tradeoff}. Given the cost ratio
\begin{equation}
\lambda = \frac{C_{\mathrm{cut}}}{C_{\mathrm{lat}}} ,
\qquad
\text{choose the policy minimising}
\quad
\lambda \cdot (\text{cut-in rate}) + (\text{latency}) .
\label{eq:cost}
\end{equation}


The units are easily lost. The cut-in rate is a fraction and the latency is in seconds, so
$\lambda$ is itself measured in seconds: $\lambda = 10$ says one interruption costs as much
as ten seconds of extra waiting, not that it is ten times worse than a delay. 

The operating point is then a consequence of a stated preference rather than a number fixed in advance.
Equation~\ref{eq:cost} is the only form of the rule used here: the policy minimising the
combined cost is selected directly, with no closed-form threshold derived from $\lambda$. A richer variant would fold the elapsed silence into the probability and
decide sequentially, in the spirit of \citet{wald1945sequential}. I built it and it did not
improve on the swept frontier (Section~\ref{sec:negative}), so the calibrated single-instant rule stands, and the full sequential formulation is left to Chapter~\ref{ch:conclusions}.
